// Convolve — blur an image using a second image as the kernel
// The @kernel image defines the blur shape (stencil).
// Useful for custom bokeh effects: use a bokeh shape image as kernel.
// Normalizes by the kernel sum so brightness is preserved.
//
// Performance note: kernel radius is limited by the loop.
// A 21x21 kernel (radius 10) is reasonable; larger gets slow.
// Maximum: radius 15 (961 iterations). Radius 16+ exceeds the 1024 iteration limit.

i$radius = 5;   // kernel half-size in pixels (5 = 11x11 stencil)

vec3 sum = vec3(0.0);
float weight_sum = 0.0;
int size = $radius * 2 + 1;
float inv_size = 1.0 / float(size - 1);
float r_offset = float($radius);

for (int ky = -$radius; ky <= $radius; ky++) {
    for (int kx = -$radius; kx <= $radius; kx++) {
        // Sample the kernel at this offset (use UV-space mapping)
        float ku = (float(kx) + r_offset) * inv_size;
        float kv = (float(ky) + r_offset) * inv_size;
        vec3 k = sample(@kernel, ku, kv).rgb;

        // Kernel weight is the luminance of the kernel pixel
        float w = luma(k);

        // Sample the source image at this offset
        vec3 s = fetch(@image, ix + kx, iy + ky).rgb;
        sum = sum + s * w;
        weight_sum = weight_sum + w;
    }
}

@OUT = sum / max(weight_sum, 0.0001);
